%
% Adaptation of the Thesis Template bei Philip Döbler shared under CC-BY 3.0
% Ergänzungen von:
% Till Tantau (teilweise entfernt)
% André Calero Valdez
% Tim Schrills


\documentclass[
    paper = a4,
    fontsize = 12pt,
    headinclude = true,
    % start chapters on the right side
    open = right,
    % Use twosided layout? Every chapter starts on the right side, therefore sometimes blank pages are added between chapters.
    twoside = false,
    BCOR = 10mm,
    % add lists and bibliography to table of contents
    toc = listofnumbered,
    toc = bibnumbered,
    % enumerate chapters as x.y instead of x.y.
    numbers = noendperiod
]{scrreprt}

% use some additional package, add your own below if required
\usepackage[utf8]{inputenc}
\usepackage[T1]{fontenc}
% package for code listings
\usepackage{listings}
% package providing colors (e.g. for syntax highlighting in listings)
\usepackage{xcolor}
% package for german language, change to english if required
\usepackage[ngerman]{babel}
\usepackage[german=quotes]{csquotes}
\usepackage{xpatch}
\usepackage{hyphenat}
% package to set the line spacing
\usepackage{setspace}
% package for hyperlinks (e.g., in the table of contents) and set links to same style as text
\usepackage[hidelinks]{hyperref}
% package to create acronyms
\usepackage[acronym]{glossaries}
% package for pseudocode
\usepackage[german, algochapter, linesnumbered]{algorithm2e}
% package to create diagrams
\usepackage{tikz}
\usepackage{booktabs}
\usepackage{float} % in der Präambel
\usepackage{lmodern}
\usepackage{array}
\usetikzlibrary{positioning}
\usetikzlibrary{shapes.geometric, arrows}
\tikzset{font={\fontsize{10pt}{12}\selectfont}}

%%% additions by ACV
% Citations using APA style with biblatex
\usepackage[style=verbose,backend=biber,backref=true,maxbibnames=20]{biblatex}
% Biblatex allows multiple bib-files
\addbibresource{bibliography/references.bib}
\addbibresource{bibliography/websources.bib}

% Tables using APA Style
\usepackage{booktabs}

% use microkerning for improved typesetting
\usepackage{microtype}

% include acronyms
% add acronyms below
\newacronym{nic}{NIC}{Network Interface Controller}


% avoid bold caption for algorithms to be consistent with other captions
\SetAlCapSty{}
\SetNlSty{footnotesize}{}{}

% definitions for syntax highlighting in listings environment
% define custom colors for syntax highlighting in code listings
\definecolor{lst_comments}{rgb}{0, 0.75, 0}
\definecolor{lst_linenumbers}{rgb}{0.2, 0.2, 0.2}
\definecolor{lst_keywords}{rgb}{0, 0, 0.75}
\definecolor{lst_strings}{rgb}{0.75, 0, 0}
\definecolor{lst_background}{rgb}{0.97, 0.97, 0.97}

% define style for listings
\lstdefinestyle{msqc_style}{
    backgroundcolor=\color{lst_background},   
    commentstyle=\color{lst_comments},
    keywordstyle=\color{lst_keywords},
    stringstyle=\color{lst_strings},
    % style of the line numbers on the left side
    numberstyle=\tiny\color{lst_linenumbers},
    % use monospace font for code and set size to the size of footnotes
    basicstyle=\footnotesize,
    xleftmargin=2em,
    % break long lines, if this happens it might be better to manually add a line break at an appropriate place
    breaklines = true,
    % place caption below the listing
    captionpos = b,
    % don't drop spaces to fix alignment and always convert tabs to spaces
    keepspaces = true,
    % place line numbers on the left side
    numbers = left,
    % distance between line numbers and listing
    numbersep = 7pt,
    % don't show special characters for spaces
    showspaces = false,
    % don't show special characters for spaces in strings
    showstringspaces = false,
    % don't show special characters for tabs
    showtabs = false,
    % set size of tab to 4 spaces
    tabsize = 4
}
% set style for the document
\lstset{style=msqc_style}

% set line spacing to 1.5
\onehalfspacing

\begin{document}
    % use roman numbering for preamble
    \pagenumbering{Roman}
    % add title page and abstract
    \begin{titlepage}
    \begin{center}
        %% TODO: find hires figure.
        \includegraphics{figures/HFT-Logo.jpg}\\
        \normalsize
        \sffamily
        %{\color{gray}Direktor:in Prof. Dr. rer. nat. Nicole Jochems\\}
        
        \vspace{0.5cm} 
        
        \Huge
        \textbf{Betreutes Praktisches Studienprojekt}
        \vfill
        
        \normalsize
        % change accordingly for Master Thesis
        \textbf{Praktikumsbericht} \\
        \vspace{0.4cm}
        
        im Rahmen des Studiengangs\\
        \textbf{Wirtschaftsinformatik}\\
        der Hochschule für Technik Stuttgart\\

        
        \vspace{0.8cm}
        vorgelegt von: \\
        \large
        \textbf{Julian Raubald}\\ 
        (Matrikelnummer 1003812)\\
        
        \normalsize
        \vspace{0.8cm}
        Praktikumsbetrieb:\\ 
        \large
        %\textbf{iteratec GmbH}\\
        \includegraphics[width=0.5\textwidth]{figures/iteratec-Logo.png}\\
        \vspace{0.3cm}
        \normalsize
        Betreut durch:\> \\ 
        \large
        \textbf{Prof. Dr. Stefan Knauth} \\
        \normalsize
        \vspace{0.5cm}
        Stuttgart, Tag. Februar. 2024\\
    \end{center}
\end{titlepage}
    
    % print table of contents
    \tableofcontents
    
    % use arabic numbering for main part
    \cleardoublepage\pagenumbering{arabic}

    % Deutsche Paragrapheinzüge auf 0 und Parskip auf 10 pt
    % Für englische Text auskommentieren
    \setlength{\parindent}{0pt}
    \setlength{\parskip}{10pt}
    
    % create a tex file for each chapter and include it below

    %%%%%%%%%%%%%%%%%%%%%%%%%%%%%%%%
    %%%%%%%%%%%%%%%%%%%%%%%%%%%%%%%% 
    % Inhalt der Arbeit

    \chapter{Einleitung}
In diesem Bericht schildere ich die Aufgaben und Projekte, welche ich in meinem Praxissemester im Wintersemester 23/24
bei dem Unternehmen iteratec GmbH durchgeführt und begleitet habe. Dabei gehe ich auf die eingesetzten Technologien ein
und teile meine gesammelten Erfahrungen und ordne diese inhaltlich im Rahmen meines Studiums der Wirtschaftsinformatik
ein.

\section{Unternehmen iteratec GmbH}

Iteratec ist ein Unternehmen, das sich auf digitale Produktinnovation, Software- und Architekturentwicklung sowie
digitale Infrastrukturen spezialisiert hat. Als Partner für die digitale Transformation unterstützt es Unternehmen und
öffentliche Organisationen bei der Digitalisierung ihrer Prozesse, Produkte und Dienstleistungen. Dies umfasst die
Entwicklung von Innovationen, die technische Umsetzung individueller Softwarelösungen, deren Betrieb und
Weiterentwicklung sowie Schulungen in agilen Methoden. Das Unternehmen entwickelt mobile Applikationen,
Cloud-Architekturen und bietet Application Performance Management an. Es nutzt auch Technologien wie Blockchain, AI und
IoT für die Entwicklung neuer Geschäftsmodelle. Zu den Kunden gehören mittelständische Unternehmen, DAX-Konzerne und
Organisationen aus verschiedenen Branchen. 1996 in München gegründet, hat iteratec Standorte in Deutschland, Österreich
und Polen und beschäftigt rund 500 Mitarbeiter und 120 Studierende. Seit 2019 sind viele Mitarbeiter über eine
Genossenschaft am Unternehmen beteiligt.

\section{Name und Geschichte}
Das Unternehmen wurde 1996 gegründet und existiert bereits seit über 25 Jahren. Während zu dieser Zeit das
Wasserfallmodell in der Softwareentwicklung weit verbreitet war hat sich iteratec schon früh auf agile Methoden und
Werte spezialisiert. Erkennbar ist das insbesondere an dem Namen, welcher sich aus den Begriffen iterativ und
Technologie zusammensetzt. Heute hat iteratec bereits über 1000 Softwareprojekte erfolgreich abgeschlossenen.

\section{Unternehmensstruktur}
Standortübergreifend hat iteratec 4 Hauptgeschäftsführer. An jedem Standort gibt es jeweils einen
Geschäftsstellenleiter. Jeder Mitarbeiter hat eine persönliche Führungskraft (pFK), welche für die persönliche
Entwicklung vor allem im, aber auch außerhalb des Unternehmens zuständig ist. Jedes Projekt hat einen Product Owner
(PO), welcher die Interessen des Kunden im Projekt vertritt, und einen Projektleiter, welcher dafür zuständig ist die
Fristen und eingesetzten Ressourcen zu managen. Alle Mitarbeiter begegnen sich Standort- und Positionsübergreifend auf
Augenhöhe, was sich unter anderem dadurch zeigt, dass sich alle duzen. Die bisherigen Alleingesellschafter der iteratec
GmbH haben im Jahr 2018 entschieden, das Unternehmen weder innerhalb ihrer Familie zu vererben noch dieses an externe zu
verkaufen. Daraus entstand die iteratec nurdemtean eG, eine Genossenschaft bei der sich die Mitarbeiter am Unternehmen
beteiligen, mit dem Ziel, die Unternehmenskultur zu halten.

\section{Studierende bei iteratec (SLAB)}
Die Abteilung für studierende, das SLAB (Studentenlabor), ist Standort übergreifend und schließt sowohl alle
Studierenden als auch deren Betreuer mit ein. Halb-jährlich werden neue SLAB Leiter gewählt, deren Aufgabe es ist, alle
anfallenden Aufgaben an studierende zu delegieren und im engen Austausch mit dem SLAB-Verantwortlichen zu stehen. Eine
Aufgabe ist das bereitstellen einer Startbegleitung für neue Studierende. Andere Aufgaben sind z.B. das organisieren von
Workshops und Events.

\section{Meetings und Veranstaltungen}
Bei iteratec gibt es abseits der regulären Projektmeetings viele weitere Veranstaltungen an denen man teilnehmen kann.
Im folgenden Kapitel erläutere ich die wichtigsten.

\subsection{SLAB-Meeting}
Das SLAB-Meeting findet zwei mal pro Jahr statt und beschränkt sich auf den Standort. Hier wird die SLAB-Leitung
gewählt, welche immer aus zwei Personen besteht. Es wird retrospektiv geschaut und bewertet wie gut die Prozesse rund um
die Studentenbetreuung funktionieren. Es werden Verbesserungsvorschläge gesammelt und nach Nutzen und Durchführbarkeit
bewertet. Anschließend werden entsprechende Maßnahmen an verantwortliche verteilt die sicherstellen, dass diese
durchgeführt werden.

\subsection{GS-Meeting}
Das Geschäftsstellen-Meeting für die Geschäftsstelle Stuttgart, an dem alle Mitarbeiter des Standorts eingeladen sind,
findet etwa ein Mal pro Monat statt. Hier werden aktuelle Themen besprochen und Mitarbeiter können Ihre Projekte
vorstellen. Die Geschäftsführung verkündet während des Meetings wichtige Neuigkeiten, anstehende Events und präsentiert
die aktuellen Geschäftszahlen. Neue Mitarbeiter und studierende werden begrüßt und ausscheidende Mitarbeiter
verabschiedet. Zum Ende können Kollegen für die Wall of Fame in einer der drei Leitwerte von iteratec (Permanent
Exzellent, Konstruktiv unabhängig und Inspirierend mutig) nominieren und so besondere Leistungen und gute zusammenarbeit
würdigen.

\subsection{BUILD23}
Bei der BUILD handelt es sich um eine Messe, welche ausschließlich von iteratec Mitarbeitern besucht werden darf. Sie
geht über zwei Tage und fand dieses Jahr bei der Geschäftsstelle in München statt. Im Vorfeld kann jeder Mitarbeiter,
Student oder jedes Team sich ein Thema überlegen und als Vorschlag einreichen. Nachdem die Themen freigegeben wurden
plant man einen entsprechenden Messestand über drei bis sechs Zeitslots, in denen man sein Thema, Projekt oder Workshop
präsentieren kann. Hierzu gibt es an jedem Stand Ausrüstung in Form eines TVs, eines Mikrofons und mehreren Receiver
mit Kopfhörern. Ziel der Build ist es das erlangte Wissen Projekt- und Standortübergreifend auszutauschen und sich mit
Kollegen zu vernetzen oder Interessensgemeinschaften zu bilden. Im Anschluss an die Build findet die Weihnachtsfeier
statt.

%
% Generelle Hinweise:
% - Werfen Sie auch einen Blick in die Word-Vorlage, falls dort Hinweise sind, die hier nicht enthalten sind.
%-  (Dieses Dokument ist für einseitigen Druck formatiert; wenn zweiseitig gedruckt werden soll, müssen die Seitenzahlen
%   und Header entsprechend angepasst werden.) 
%-  Auf Abbildungen / Tabellen wird möglichst im Text vor der Abbildung verwiesen. Ist in Latex manchmal schwierig
%-  Abbildungen sollten nach Möglichkeit so groß dargestellt sein, dass auch die Texte gut lesbar sind; es sei denn
%   diese sind völlig bedeutungslos und nur die Struktur oder das Gesamtbild sind von Bedeutung.
%-  Bei farbigen Abbildungen sollte sichergestellt werden, dass diese auch in Schwarz-Weiß gut erkennbar sind.
%-  Tabellen sollten zweckmäßig und übersichtlich sein: Vermeidung unnötiger Linien, Farbgebung nur, wenn sie eine
%   Bedeutung hat oder der Übersichtlichkeit dient.
%-  Zitiert wird typischerweise nach APA. Alternativen sind aber möglich (mit den Betreuer:innen klären). 


    \chapter{Grundlagen der KI und des Datenschutzes}

\section{Künstliche Intelligenz (KI)}

Künstliche Intelligenz (KI) bezeichnet Technologien, die es Computern
ermöglichen, Aufgaben welche normalerweise menschliche Intelligenz benötigen.
Wichtige Bestandteile sind dabei das Lernen, Problemlösen und verstehen von
Sprache. Ein zentrales Element ist dabei das maschinelle Lernen, bei dem riesige
Datenmengen mithilfe von Algorithmen Muster formen und Vorhersagen treffen. 
Mit der Verbreitung von Large Language Models (LLMs) wie GPT-3.5 von OpenAI hat
die Zugänglichkeit und Nutzung von KI stark zugenommen. 
\cite{haerting2024}


\section{Datenschutz}

Datenschutz schützt personenbezogene Daten vor Missbrauch und unbefugtem
Zugriff. In der EU regelt die Datenschutz-Grundverordnung (DSGVO) diesen Schutz.
Personenbezogene Daten umfassen sämtliche Informationen, die eine Person identifizieren,
wie Namen, Adressen und pseudonymisierte Daten wie Kundennummern. Die DSGVO erlaubt
die Verarbeitung solcher Daten nur auf rechtmäßiger Grundlage, wie
durch eine Einwilligung oder das berechtigte Interesse des Datenverarbeiters.

\cite{haerting2024}

    \chapter{Herausforderungen und Risiken}

In diesem Kapitel werden Beispiele für die Herausforderungen und Risiken im Zusammenhang mit
der Nutzung von KI im Bereich des Datenschutzes erläutert. 

% Dabei liegt der Fokus auf den potenziellen Datenschutzbeeinträchtigungen durch KI.

\section{Datenlecks}

KI-Systeme können anfällig für Datenlecks sein, bei denen sensible Informationen
unbefugt offengelegt werden. Ein Datenleck kann durch Schwachstellen in der
IT-Infrastruktur oder durch menschliches Versagen entstehen. Solche Vorfälle
können schwerwiegende Folgen haben, wie Identitätsdiebstahl oder finanziellen
Betrug. Um das Risiko von Datenlecks zu minimieren, müssen Unternehmen
sicherstellen, dass sie über robuste Sicherheitsmaßnahmen verfügen, wie z.B.
regelmäßige Sicherheitsüberprüfungen und den Einsatz von
Verschlüsselungstechnologien.

\section{Profiling und Diskriminierung}

Algorithmen können personenbezogene Daten analysieren und daraus Profile
erstellen, die zu Diskriminierung führen können. Dies geschieht häufig unbemerkt
und kann tiefgreifende Auswirkungen auf die betroffenen Personen haben.
Beispielsweise könnten Bewerber aufgrund von algorithmischen Entscheidungen von
Bewerbungsverfahren ausgeschlossen werden oder Verbraucher könnten unfairen
Kreditentscheidungen ausgesetzt sein. Unternehmen müssen sicherstellen, dass
ihre KI-Systeme fair und transparent sind und regelmäßig auf
Diskriminierungspotenzial überprüft werden.

\section{Unerwünschte Datenverarbeitung}

KI-Systeme könnten personenbezogene Daten für Zwecke nutzen, die nicht im
Einklang mit den ursprünglichen Erhebungszwecken stehen. Dies kann zu einer
unerwünschten Verarbeitung von Daten führen, die gegen Datenschutzbestimmungen
verstößt. Beispielsweise könnten Daten, die für die Verbesserung eines Dienstes
erhoben wurden, ohne Zustimmung der betroffenen Personen für Marketingzwecke
verwendet werden. Unternehmen müssen klare Richtlinien für die Datennutzung
festlegen und sicherstellen, dass die Verwendung von Daten stets im Einklang mit
den ursprünglichen Erhebungszwecken und den Einwilligungen der Nutzer steht.

    % \chapter{Gesetzliche Rahmenbedingungen}

\section{Datenschutzgesetz A}

\section{Datenschutzgesetz B}
    \chapter{Lösungsansätze}

\section{Technische Maßnahmen}
Ein zentraler Lösungsansatz zur Bewältigung der Datenschutzprobleme im
Zusammenhang mit Künstlicher Intelligenz (KI) besteht in der Implementierung
technischer Maßnahmen. Eine wichtige Maßnahme ist die Verschlüsselung von Daten,
um sensible Informationen vor unbefugtem Zugriff zu schützen. Unternehmen
sollten moderne Verschlüsselungstechnologien einsetzen und regelmäßig
aktualisieren.

Anonymisierungs- und Pseudonymisierungstechniken sind ebenfalls entscheidend.
Anonymisierte Daten unterliegen nicht den strengen Vorgaben der
Datenschutz-Grundverordnung (DSGVO), und pseudonymisierte Daten können das
Datenschutzrisiko verringern \cite{conrad_2017}.

Regelmäßige Sicherheitsüberprüfungen und Penetrationstests sind notwendig, um
potenzielle Schwachstellen in IT-Systemen zu identifizieren und zu beheben.
Diese Maßnahmen minimieren das Risiko von Datenlecks und unbefugtem Zugriff.

\section{Organisatorische Maßnahmen}
Neben technischen Maßnahmen spielen auch organisatorische Maßnahmen eine
wichtige Rolle beim Schutz personenbezogener Daten. Dazu gehört die Entwicklung
klarer Datenschutzrichtlinien und -verfahren sowie die Schulung von Mitarbeitern
im Umgang mit sensiblen Daten.

Datenschutz-Folgenabschätzungen (DSFA) sind gemäß Art. 35 DSGVO erforderlich,
bevor neue Technologien eingeführt werden. Diese Bewertungen helfen, potenzielle
Datenschutzrisiken frühzeitig zu erkennen und geeignete Maßnahmen zu ergreifen
\cite{conrad_2017}.

Die Ernennung eines Datenschutzbeauftragten, der die Einhaltung der
Datenschutzvorschriften überwacht, ist ebenfalls wichtig. Der
Datenschutzbeauftragte sollte Schulungen und Workshops für Mitarbeiter
organisieren, um das Bewusstsein für Datenschutzthemen zu schärfen.

Klare Richtlinien für die Nutzung und Verarbeitung personenbezogener Daten sind
entscheidend. Diese Richtlinien sollten die Zwecke der Datenerhebung und
-verarbeitung genau definieren und sicherstellen, dass personenbezogene Daten
nicht für andere Zwecke verwendet werden, als ursprünglich vorgesehen.
Mechanismen zur Wahrung der Rechte der betroffenen Personen auf Auskunft,
Berichtigung und Löschung ihrer Daten sollten implementiert werden
\cite{conrad_2017}.

Durch die Kombination technischer und organisatorischer Maßnahmen können
Unternehmen die Datenschutzrisiken im Zusammenhang mit der Nutzung von KI
erheblich reduzieren und sicherstellen, dass ihre KI-Systeme im Einklang mit den
gesetzlichen Datenschutzanforderungen stehen.


    \chapter{Abschluss und Ausblick}

\section{Zusammenfassung wichtigster Punkte}
In diesem Bericht wurde detailliert untersucht, wie Künstliche Intelligenz (KI)
sowohl Chancen als auch Herausforderungen im Bereich des Datenschutzes mit sich
bringt. Zunächst wurden die Grundlagen der KI und des Datenschutzes erläutert.
Dabei wurde hervorgehoben, dass KI-Systeme in der Lage sind, große Datenmengen
zu verarbeiten und daraus wertvolle Erkenntnisse zu gewinnen. Gleichzeitig
wurden die datenschutzrechtlichen Anforderungen der DSGVO betont, die
sicherstellen sollen, dass personenbezogene Daten nur rechtmäßig verarbeitet
werden dürfen.

Im Kapitel über Herausforderungen und Risiken wurden spezifische Probleme wie
Profiling und Diskriminierung sowie unerwünschte Datenverarbeitung beleuchtet.
Diese Risiken verdeutlichen, dass die Nutzung von KI-Systemen strenge
Datenschutzvorkehrungen erfordert, um die Rechte der betroffenen Personen zu
schützen.

Die Lösungsansätze konzentrierten sich auf technische und organisatorische
Maßnahmen. Technische Maßnahmen umfassen Verschlüsselung, Anonymisierung und
regelmäßige Sicherheitsüberprüfungen. Organisatorische Maßnahmen betonen die
Bedeutung von Datenschutz-Folgenabschätzungen, klaren Richtlinien und der
Schulung von Mitarbeitern.

\section{Ausblick auf die Zukunft von KI und Datenschutz}
Der Einsatz von Künstlicher Intelligenz wird in den kommenden Jahren
voraussichtlich weiter zunehmen und tiefgreifende Veränderungen in vielen
Lebensbereichen bewirken. Während KI-Technologien weiterentwickelt werden,
müssen die Datenschutzregelungen entsprechend angepasst werden, um neuen
Herausforderungen gerecht zu werden.

Es ist zu erwarten, dass zukünftige Entwicklungen im Bereich der KI strengere
gesetzliche Vorgaben und fortschrittlichere technische Lösungen erfordern
werden, um den Datenschutz zu gewährleisten. Ein wichtiger Aspekt wird die
Entwicklung von KI-Systemen sein, die von Anfang an datenschutzfreundlich
konzipiert sind (Privacy by Design und Privacy by Default).

Zusätzlich könnte die internationale Zusammenarbeit bei der Entwicklung von
Datenschutzstandards und -richtlinien intensiviert werden, um einen globalen
Schutz der Privatsphäre zu gewährleisten. Forschung und Innovation im Bereich
der Datenschutztechnologien werden weiterhin eine zentrale Rolle spielen, um mit
den schnellen Fortschritten der KI Schritt zu halten.

Insgesamt zeigt sich, dass der verantwortungsbewusste Einsatz von Künstlicher
Intelligenz nur möglich ist, wenn Datenschutz und ethische Überlegungen fest in
den Entwicklungs- und Implementierungsprozessen verankert sind. Unternehmen,
Regierungen und die Gesellschaft als Ganzes müssen zusammenarbeiten, um eine
Balance zwischen technologischem Fortschritt und dem Schutz der individuellen
Privatsphäre zu finden.


    % add bibliography
    %\printbibheading
    % \printbibliography[nottype=online,notkeyword=norm]
    \printbibliography[
    title={Quellen}]
    % \printbibliography[keyword=norm,heading=subbibliography,
    % title={Normen}]

    
    % \begin{appendix}
    %     % create a tex file for each appendix and include it below
    %     \chapter{Anhang}
Der Anhang kann Teile der Arbeit enthalten, die im Hauptteil zu weit führen würden, aber trotzdem für manche Leser interessant sein könnten. Das können z.\,B. die Ergebnisse weiterer Messungen sein, die im Hauptteil nicht betrachtet werden aber trotzdem durchgeführt wurden. Es ist ebenfalls möglich längere Codeabschnitte anzuhängen. Jedoch sollte der Anhang kein Ersatz für ein Repository sein und nicht einfach den gesamten Code enthalten.
        
    %     % add lists of figures, tables and listings to the appendix
    %     \listoffigures
    %     \listoftables
    %     \lstlistoflistings
        
    %     % add acronyms
    %     \printglossary[type=\acronymtype, title=Abkürzungsverzeichnis, toctitle=Abkürzungsverzeichnis, numberedsection] 
        

    % \end{appendix}
    % \newpage
    % \Large
Erklärung
\newline
\vspace{3\baselineskip}
\normalsize
\noindent

Ich versichere an Eides statt, die vorliegende Arbeit selbstständig verfasst und nur die angegebenen Quellen benutzt zu haben.

%%%%
% HINWEIS: Was bedeutet eine eidesstattliche Versicherung?
%% 
%Die Strafbarkeit einer falschen eidesstattlichen Versicherung ist mir bekannt, namentlich die Strafandrohung gemäß § 156 StGB bis zu drei Jahren Freiheitsstrafe oder Geldstrafe bei vorsätzlicher Begehung der Tat bzw. gemäß § 163 Abs.1 StGB bis zu einem Jahr Freiheitsstrafe oder Geldstrafe bei fahrlässiger Begehung. 

\vspace{5\baselineskip}


Unterschrift 

Lübeck, Tagesdatum


    %%%%%%%%%%%%%%%%%%%%%%%%%%%%%%%%
    %%%%%%%%%%%%%%%%%%%%%%%%%%%%%%%% 
    
\end{document}
